%% Case Studies — \neven{} Paper (Revised with TipoOutput explanation + German Credit)

\section{Case Studies}
\label{sec:case_studies}

Before presenting the validation scenarios, we describe the \texttt{TipoOutput} parameter---a central design element that enables users to extract multiple perspectives from a single analytical function without writing any code.

\subsection{The TipoOutput Paradigm}
\label{subsec:tipooutput}

Every analytical function in \neven{} accepts a \texttt{TipoOutput} parameter that controls which result is returned to Excel. This design allows a single function to serve multiple analytical purposes:

\begin{itemize}
    \item \texttt{TipoOutput = 0} --- Returns the list of all available procedures for that function. This serves as built-in documentation: the user can discover what outputs are available without consulting external references.
    \item \texttt{TipoOutput = 1, (...) , N-1} --- Returns a specific result. For example, in \texttt{=R.MR\_Lineal(Y, X, TipoOutput)}: type~1 returns the estimated model with standard errors and significance; type~2 returns in-sample predictions; type~3 returns out-of-sample predictions; type~4 returns marginal effects; and so on.
    \item \texttt{TipoOutput = N} --- Returns the \emph{Universal Output Extraction}: a complete structured table containing all model outputs (coefficients, test statistics, residuals, fit measures) in a single multi-cell array. This is powered by the \texttt{Extraer\_outputs} function, which decomposes any R model object into a standardized [Section, Parameter, Metric, Value] format.
\end{itemize}

This paradigm eliminates the need for users to understand the internal structure of R model objects. A researcher exploring a regression model can iterate through \texttt{TipoOutput} values to examine different facets of the same estimation, or use \texttt{TipoOutput = 0} to discover what is available. The Universal Output Extraction (\texttt{TipoOutput = N}) provides a comprehensive view suitable for reporting or further analysis within Excel.

\paragraph{The Observation Filter (FILTRO).}
In addition to \texttt{TipoOutput}, every analytical function includes a \texttt{FILTRO} parameter: a binary vector of the same length as the dataset where a value of~1 excludes the corresponding observation from the computation. This mechanism allows the user to mark specific rows directly in Excel and instantly recalculate the model without those observations---quantifying the effect of individual data points on the results.

The pedagogical value of this design is significant. In traditional tools, excluding observations requires manipulating the original dataset (deleting rows, creating filtered copies, or writing conditional logic). In \neven{}, the user simply changes a~0 to a~1 in a cell and the entire model recalculates. This enables sensitivity analysis at the observation level: a student can identify an outlier, exclude it, and immediately observe how coefficients, standard errors, and fit statistics change. The tool becomes not merely a computation engine but an instrument for deep understanding of how each data point influences the analytical outcome.

For example, in a regression with $n = 526$ observations, the formula:
\begin{center}
\texttt{=R.MR\_Lineal(Y, X, FILTRO, 0, 0, , 1, 1)}
\end{center}
where \texttt{FILTRO} is a column of zeros and ones, allows the researcher to iteratively exclude suspected influential observations and compare the resulting models---all within the same worksheet, with no code and no data manipulation.

\paragraph{AI-Assisted Interpretation.}
Every analytical function in \neven{} can be complemented with AI-generated interpretation through the Python integration. The function \texttt{=P.ai\_call(data, "prompt")} sends the output of any statistical procedure to a large language model---either cloud-based (OpenAI, Azure OpenAI) or locally hosted (Ollama, LM Studio)---and returns a natural-language explanation directly in Excel cells. This capability is transversal: it applies to any function output regardless of the source language (R, Julia, or Python). The user configures the AI provider connection (API endpoint, model name, and authentication key) in \texttt{neven-config.json}, requiring no code changes to switch between providers or models.

The interpretation prompts are fully customizable. \neven{} ships with a set of default prompts tailored to each analytical domain (regression diagnostics, time series forecasting, clustering interpretation), but the user can modify, extend, or replace them by editing simple text files in the configuration directory. This means that a professor can craft pedagogically appropriate prompts for students, while a risk analyst can write prompts aligned with regulatory language. The AI layer transforms raw statistical output into actionable insight without requiring the user to interpret $p$-values, eigenvalues, or residual patterns manually---while preserving full transparency, since the underlying numerical results remain visible in adjacent cells.

All examples in this section reference Excel workbooks distributed with \neven{} in the \texttt{Ejemplos/Excel/} directory. Readers can open these workbooks, modify the data or parameters, and verify that the formulas produce the reported results---ensuring full reproducibility of every claim made in this paper.

\subsection{Econometric Validation}
\label{subsec:econometric}

To validate the correctness of \neven{}'s statistical implementations, we reproduce selected analyses from \citet{Wooldridge2016}, a standard reference in introductory econometrics used in graduate programs worldwide. The accompanying Excel workbooks (\texttt{CAP03.xlsm} through \texttt{CAP19.xlsm} in the \texttt{Ejemplos/Excel/WOOLDRIDGE/} directory) contain the original datasets and NEVEN formulas for each chapter. For each model, we compare three execution paths: direct R code in RStudio, the equivalent NEVEN Excel formula, and published textbook results.

\paragraph{OLS Regression (Chapter 6).}
The wage equation from Example~6.1 models $\log(\text{wage}) = \beta_0 + \beta_1 \text{educ} + \beta_2 \text{exper} + \beta_3 \text{tenure} + u$ using the \texttt{wage1} dataset ($n = 526$). In \neven{}, the user enters:

\begin{center}
\texttt{=R.MR\_Lineal(B2:B527, C2:E527, 0, 0, 0, , 1, 1)}
\end{center}

where column~B contains the dependent variable, columns~C--E contain regressors, and the final parameter \texttt{1} specifies \texttt{TipoOutput = 1} (coefficient table with standard errors, $t$-statistics, and $p$-values). Table~\ref{tab:ols_results} presents the comparison across all three execution paths.

\begin{table}[htbp]
\centering
\caption{OLS regression results for the wage equation (Chapter~6, Example~6.1). Data: \texttt{wage1} dataset ($n = 526$). Source: \texttt{CAP06.xlsm}. All three methods produce identical estimates to four decimal places.}
\label{tab:ols_results}
\begin{tabular}{@{}lcccc@{}}
\toprule
\textbf{Variable} & \textbf{Wooldridge (2016)} & \textbf{Direct R} & \textbf{\neven{}} & \textbf{Difference} \\
\midrule
Intercept       & 0.2844  & 0.2844  & 0.2844  & 0.0000 \\
Education       & 0.0920  & 0.0920  & 0.0920  & 0.0000 \\
Experience      & 0.0041  & 0.0041  & 0.0041  & 0.0000 \\
Tenure          & 0.0221  & 0.0221  & 0.0221  & 0.0000 \\
\addlinespace
$R^2$           & 0.316   & 0.3160  & 0.3160  & 0.0000 \\
$\bar{R}^2$     & 0.312   & 0.3121  & 0.3121  & 0.0000 \\
$F$-statistic   & 80.39   & 80.39   & 80.39   & 0.00   \\
\bottomrule
\end{tabular}
\end{table}

Using \texttt{TipoOutput = 13} on the same formula returns the Universal Output Extraction---a complete table with all model components (coefficients, residuals, $R^2$, $F$-statistic, AIC, BIC) in a single multi-cell array that the user can filter and analyze within Excel.

\paragraph{Binary Response Models (Chapter 17).}
For the probit model of married women's labor force participation, $P(\text{inlf} = 1 | \mathbf{x}) = \Phi(\beta_0 + \beta_1 \text{nwifeinc} + \beta_2 \text{educ} + \beta_3 \text{exper} + \beta_4 \text{exper}^2)$, the \neven{} formula \texttt{=R.MR\_Binario(Y, X, 0, 0, 0, , 1, 1)} produces marginal effects and goodness-of-fit statistics consistent with the textbook values. The Hosmer--Lemeshow test statistic and classification accuracy match those reported by \citet{Wooldridge2016} within rounding precision. Source workbook: \texttt{CAP17.xlsm}.

\paragraph{Panel Data (Chapter 13).}
Panel data models require the \texttt{plm} package \citep{Croissant2008} in R. \neven{} exposes this functionality through \texttt{=R.MR\_PanelData(Y, X, ID, Time, TipoOutput)}, where \texttt{ID} identifies the cross-sectional unit and \texttt{Time} identifies the time period. The \texttt{TipoOutput} parameter selects between coefficient estimates~(1), entity fixed effects~(2), Hausman test~(4), and complete output extraction~(16). Table~\ref{tab:panel_results} compares the fixed effects estimates. Source workbook: \texttt{CAP13.xlsm}.

\begin{table}[htbp]
\centering
\caption{Panel data fixed effects estimates for county crime rates (Chapter~13). Data: \texttt{crime2} dataset (46 cities, 7 years). Source: \texttt{CAP13.xlsm}. \neven{} results match direct R execution exactly.}
\label{tab:panel_results}
\begin{tabular}{@{}lccc@{}}
\toprule
\textbf{Statistic} & \textbf{Direct R (\texttt{plm})} & \textbf{\neven{}} & \textbf{Match} \\
\midrule
Unemployment coeff.  & 0.0342  & 0.0342  & \checkmark \\
Robust SE            & 0.0117  & 0.0117  & \checkmark \\
Hausman $\chi^2$     & 14.73   & 14.73   & \checkmark \\
Hausman $p$-value    & 0.0001  & 0.0001  & \checkmark \\
Within $R^2$         & 0.1282  & 0.1282  & \checkmark \\
\bottomrule
\end{tabular}
\end{table}

\paragraph{Time Series (Chapter 11)}
The time series module (\small \texttt{=R.ST\_SeriesTemporales}) implements unit root tests, ARIMA, GARCH, and decomposition models. For the Phillips curve ($\Delta \text{inf}_t = \alpha_0 + \alpha_1 \text{unem}_t + e_t$), \neven{} produces identical AIC/BIC values, coefficient estimates, and Ljung--Box test statistics compared to direct R execution using the \texttt{forecast} package \citep{Hyndman2008}. Source workbook: \texttt{CAP11.xlsm}.

\paragraph{Julia Numerical Validation.}
The Julia module provides 70 procedures across 10 modules. We validated the linear algebra module (\texttt{=J.Algebra(...)}) against known matrix decomposition results. Table~\ref{tab:julia_validation} confirms that all decompositions achieve machine-precision accuracy, demonstrating that Protocol Buffers serialization and Named Pipe transport do not introduce numerical degradation.

\begin{table}[htbp]
\centering
\caption{Julia numerical validation: matrix decomposition accuracy for a $100 \times 100$ random matrix. Errors measured as $\|\mathbf{A} - \hat{\mathbf{A}}\|_F / \|\mathbf{A}\|_F$.}
\label{tab:julia_validation}
\begin{tabular}{@{}lcc@{}}
\toprule
\textbf{Decomposition} & \textbf{\neven{} Formula} & \textbf{Relative Error} \\
\midrule
LU factorization    & \texttt{=J.Algebra(A, , "lu")}  & $< 10^{-15}$ \\
QR factorization    & \texttt{=J.Algebra(A, , "qr")}  & $< 10^{-15}$ \\
SVD                 & \texttt{=J.Algebra(A, , "svd")} & $< 10^{-15}$ \\
Eigendecomposition  & \texttt{=J.Algebra(A, , "eigen")} & $< 10^{-14}$ \\
Cholesky (SPD)      & \texttt{=J.Algebra(A, , "cholesky")} & $< 10^{-15}$ \\
\bottomrule
\end{tabular}
\end{table}

\paragraph{Reproducibility}
The Excel workbooks in \texttt{Ejemplos/Excel/WOOLDRIDGE/} (chapters 3--19) and \texttt{Ejemplos/Excel/BONILLA/} (specialized procedures) are distributed with the software. The complete validation workflow is fully reproducible. A researcher can open any workbook, verify that the formulas produce the expected results, modify the data or model specification interactively, and use \texttt{TipoOutput = 0} to discover additional analytical options. Combined with Quarto integration (\texttt{=NEVEN.q("validation\_report.qmd")}), the entire analysis pipeline from raw data to formatted report is auditable.

\paragraph{Aerospace Engineering: HL-20 Reentry Trajectory.}
To demonstrate \neven{}'s applicability beyond statistics and economics, we include a simplified atmospheric reentry simulation inspired by the NASA HL-20 lifting body vehicle. The 3-DOF model (altitude, velocity, flight path angle) is implemented as a Pluto.jl reactive notebook (\texttt{hl20\_reentry.jl}) that receives initial conditions from Excel via \texttt{=NEVEN.pluto.data(A1:C1, "hl20\_params")} and solves the trajectory using a Runge-Kutta~4 integrator. The user defines entry altitude ($h_0 = 120$~km), orbital velocity ($V_0 = 7500$~m/s), and flight path angle ($\gamma_0 = -1.5^\circ$) in Excel cells, sends them to Julia with a single formula, and observes the resulting trajectory---peak dynamic pressure, maximum g-load, descent profile---interactively in the Pluto notebook.

This example illustrates a key design principle: \neven{} does not impose ceilings on computational complexity. The same infrastructure that serves an economics student running a linear regression also serves an aerospace engineer solving coupled differential equations. Each user extends the platform to their level of rigor---from the simplified 3-DOF model demonstrated here to a full 6-DOF simulation with Mach-dependent aerodynamic tables (NASA~TM-4302), thermal constraints, and guidance algorithms---simply by adding Julia code to the notebook. The Excel-to-Julia pipeline handles data flow; the scientific complexity is limited only by the user's domain expertise.

\subsection{Enterprise Analytics: Credit Risk with German Credit Data}
\label{subsec:enterprise}

To demonstrate \neven{}'s applicability in enterprise settings, we present a credit risk analysis using the German Credit dataset \citep{Dua2019}---a publicly available benchmark containing 1,000 loan applications with 20 predictor variables and a binary outcome (good/bad credit risk). The accompanying workbook (\texttt{Ejemplos/Excel/BONILLA/R4XL-GermanCredit.xlsx}) allows readers to reproduce every result presented in this section.

\paragraph{Data Acquisition and Exploration.}
The workflow begins entirely within Excel. The user retrieves the dataset directly from R's package ecosystem using a single formula:

\begin{center}
\texttt{=NEVEN.r("data('german', package='rchallenge'); print(as.data.frame(german))")}
\end{center}

Once the data resides in Excel, the user explores each variable's distinct values with \texttt{=R.DB\_Unicos(A1:A1001)}, identifying the categorical structure of the dataset without writing any R code. Variables encoded as text (e.g., \texttt{A11}, \texttt{A12} for account status) can be recodified into human-readable labels using \texttt{=R.DB\_Recodificar(range, old\_values, new\_values)}. A training set is constructed via \texttt{=R.FX\_AleatorioUniforme(1000, 0, 1, seed)} to generate a random uniform vector that partitions observations into training and test subsets---all using native Excel formulas.

\paragraph{Automatic Categorical Variable Detection.}
A key innovation in \neven{} is the automatic detection of variable types during model estimation. When the user passes a dataset containing non-numeric variables (such as the German Credit's account status, credit history, or employment type), the system internally identifies which columns are categorical and automatically computes the effect of each class on the model---equivalent to the manual creation of dummy variables that traditional tools require. The user simply passes the data as-is; \neven{} handles the factor encoding transparently, producing coefficient estimates for each category level without any additional configuration.

\paragraph{Credit Scoring Model.}
The logistic regression for probability of default uses:

\begin{center}
\texttt{=R.MR\_Binario.C(Y, X, 0, Filtro, 0, 0, 0, TipoOutput)}
\end{center}

where \texttt{Y} contains the binary credit risk outcome, \texttt{X} contains all 20 predictor variables (numeric and categorical mixed), and \texttt{Filtro} defines the training subset. \texttt{TipoOutput = 1} returns coefficient estimates including the effect of each categorical level. \texttt{TipoOutput = 3} returns marginal effects at the mean. \texttt{TipoOutput = 5} returns ROC curve data for model discrimination assessment. The Universal Output Extraction (\texttt{TipoOutput = N}) returns all model components in a single structured table---enabling the analyst to build a complete model assessment without writing any R code.

\paragraph{Interactive Visualization.}
The dashboard function \texttt{=R.Dashboard(A1:U1001, 1)} generates a multi-tab interactive interface combining rpivotTable, Plotly charts, and D3.js visualizations in a single WebView2 window. The Snap Layout feature automatically positions Excel on the left and the viewer on the right, enabling iterative exploration: the analyst modifies data or parameters in Excel and immediately observes the impact in the visualization.

\paragraph{Regulatory Documentation.}
The formula \texttt{=NEVEN.q("credit\_risk\_report.qmd")} renders a Quarto document containing the analysis narrative, embedded R code chunks referencing the Excel data, and cross-references to model outputs. This produces a publication-quality report suitable for regulatory submission (Basel~III/IV compliance), with all computations reproducible from the source document.

\paragraph{Workflow Comparison.}
Table~\ref{tab:enterprise_comparison} contrasts the \neven{}-based workflow with the traditional multi-tool approach for the same credit risk analysis.

\begin{table}[htbp]
\centering
\caption{Workflow comparison: traditional enterprise analytics vs.\ \neven{}-based approach for credit risk analysis on the German Credit dataset.}
\label{tab:enterprise_comparison}
\begin{tabular}{@{}lp{4.5cm}p{4.5cm}@{}}
\toprule
\textbf{Aspect} & \textbf{Traditional} & \textbf{\neven{}} \\
\midrule
Tools required & Excel + SAS/SPSS + Tableau + Word & Excel + \neven{} (single environment) \\
\addlinespace
Skills required & VBA, SAS programming, Tableau & Excel formulas only \\
\addlinespace
Handoff points & 4 (data export/import between tools) & 0 (all in Excel) \\
\addlinespace
Reproducibility & Manual (re-run each tool) & Automatic (recalculate formulas) \\
\addlinespace
Licensing & Commercial (restricted access) & Open source (GPL v3) \\
\addlinespace
Audit trail & Fragmented across tools & Unified in Excel + Quarto \\
\bottomrule
\end{tabular}
\end{table}

The elimination of handoff points reduces both analysis time and the risk of data transformation errors. The open-source nature of \neven{} removes licensing barriers that often exclude academic institutions and small enterprises from accessing advanced analytical capabilities.
